\chapter{Elementos de Diseño}

Cada elemento dentro del diseño conceptual posee propiedades específicas que determinan su comportamiento sonoro y funcional. A continuación se detallan los elementos disponibles organizados por su función principal.

\section{Generadores de Sonido}

Los generadores son los encargados de producir la señal de audio inicial.

\subsection{Oscilador (Oscillator)}
Genera formas de onda periódicas (senoidal, cuadrada, sierra, etc.) con control de frecuencia, fase y ancho de pulso.

\subsection{Ruido (Noise)}
Genera señales aleatorias de diferentes tipos (blanco, rosa, marrón) útiles para efectos de percusión o texturas atmosféricas.

\subsection{Sampler (Sample)}
Permite la reproducción de archivos de audio externos, con opciones de bucle (loop) y control de ganancia.

\section{Modificadores de Parámetros}

Estos elementos no generan audio por sí mismos, sino que producen señales de control para modificar parámetros de otros nodos.

\subsection{LFO (Low Frequency Oscillator)}
Oscilador de baja frecuencia utilizado para modulaciones cíclicas (vibrato, trémolo, etc.).

\subsection{Envolvente (Envelope)}
Generador de señales ADSR (Attack, Decay, Sustain, Release) para controlar la evolución temporal de parámetros como el volumen o el corte de filtro.

\section{Modificadores de Sonido y Efectos}

Procesan la señal de audio entrante para alterar su timbre o dinámica.

\subsection{Filtro de Frecuencia (Frequency Filter)}
Permite el paso de ciertas frecuencias mientras atenúa otras (paso bajo, paso alto, paso banda).

\subsection{Reverberación (Reverb)}
Simula la acústica de un espacio físico, añadiendo profundidad y persistencia al sonido.

\subsection{Retraso (Delay)}
Repite la señal de audio tras un intervalo de tiempo definido, creando ecos.

\subsection{Distorsión (Distortion)}
Altera la forma de onda de la señal para añadir armónicos y saturación.

\subsection{Chorus / Flanger}
Efectos de modulación basados en el tiempo que añaden grosor y movimiento espacial.

\subsection{Compresor (Compressor)}
Controla el rango dinámico de la señal, reduciendo el nivel de los picos más altos.

\subsection{Ecualizador (Equalizer)}
Permite ajustar selectivamente la ganancia de diferentes bandas de frecuencia.

\section{Mezcla y Control Espacial}

Elementos destinados a la gestión de múltiples señales y su posicionamiento.

\subsection{Mezclador (Mixer)}
Combina hasta 10 señales de entrada en una o dos salidas estéreo.

\subsection{Ganancia y Panoramización (Gain and Pan)}
Control básico de nivel de salida y posicionamiento en el campo estéreo.
